The network gives us a way to describe how characteristic a recipe's ingredient pairs are of its Food.com cuisine tag. It does not tell us whether the dish is culturally authentic. The next step is to compare this score with what reviewers actually say, distinguishing comments that affirm authenticity from those that question it. A review with no mention of authenticity cannot be read as a negative judgment. General sentiment and star ratings should be examined separately.

The two randomizations answer different questions. Shuffling cuisine tags shows that the recipe network is associated with those tags. Preserving ingredient frequencies shows that much of this association can arise because cuisines use different individual ingredients. The smaller remaining difference cannot be attributed uniquely to cuisine-specific pairing rules: recipe families and preparation styles may also contribute. The pair check makes another distinction concrete: olive oil and parmesan are distinctive of Italian-tagged recipes relative to other cuisines, yet occur together less often than the frequency-preserving reference predicts within Italian-tagged recipes. Our recipe pair score measures cross-cuisine prevalence, not a pairing effect independent of ingredient popularity. We will therefore ask whether pair scores add information beyond individual ingredients, while comparing recipes within cuisines and accounting for their review counts and broad dish types. A null result would mean these measurable signatures do not predict how this Food.com audience uses the word ``authentic''; it would not settle the meaning of authenticity more broadly.

Food.com cuisine and course tags are supplied by users, and our manual review found uncertain labels. The ingredient cleaning deliberately merges only clear aliases, leaving some equivalent names separate. Preparation methods, personal memories, and access to ingredients may shape a review but are absent from the pair score. The dataset represents an online recipe community, not all cooking within any of the named cultures. These limits affect both the score and its interpretation alongside review text.
